% Slide 1 – title. Team name, ID and links are set in main-en.tex.
{
\setbeamercolor{background canvas}{bg=accNavy}
\setbeamercolor{normal text}{fg=white}
\usebeamercolor[fg]{normal text}
\hypersetup{urlcolor=white}   % links on the navy background; outside this group they return to the main.tex colour
\begin{frame}[plain]
  \vspace{6mm}
  \begin{columns}[T,onlytextwidth]
    \begin{column}{0.64\textwidth}
      \begin{tikzpicture}
        \node[fill=white,rounded corners=2.5mm,inner sep=1mm] {\includegraphics[height=12mm]{logo}};
      \end{tikzpicture}

      \vspace{2mm}
      {\usebeamerfont{title}\color{white}Accessly\par}
      \vspace{0.5mm}
      {\Large\bfseries\color{white}Kraków without barriers\par}
      \vspace{1.5mm}
      {\color{accSignal}Your city. Your needs. Check before you go.\par}
      \vspace{3mm}
      {\small\color{white}
        A guide to accessible places and a map of barriers in public space.
        Every piece of information carries its source, date and verification status.\par
        \vspace{2.5mm}
        \textbf{\TeamName}\ \textperiodcentered\ \TeamId\par
        HackYeah 2026 \textperiodcentered\ “Cracow without barriers” challenge\par
        \vspace{2.5mm}
        {\footnotesize
        Demo: \DemoUrl\\
        Code: \RepoUrl\\
        Video (3~min): \VideoUrl\par}}
    \end{column}
    \begin{column}{0.32\textwidth}
      \centering
      \vspace{-1mm}
      \includegraphics[height=70mm]{mobile-mapa}
    \end{column}
  \end{columns}

  \note[item]{One sentence: Accessly does not say “this place is accessible”; it says “this place will work for you”, and shows where we know it from, since when, and how far it can be trusted.}
  \note[item]{Show the phone next to the laptop: the same prototype runs on both.}
  \note[item]{Do not read the links out loud; they are in the PDF.}
\end{frame}
}
